% Select the legacy font encoding before IEEEtran initializes Times fonts.
% Tectonic/XeTeX otherwise defaults to TU, for which ptm has no font shapes.
\RequirePackage[T1]{fontenc}
\documentclass[conference]{IEEEtran}
\IEEEoverridecommandlockouts
\usepackage{cite}
\usepackage{amsmath,amssymb,amsfonts}
\usepackage{graphicx}
\usepackage{textcomp}
\usepackage{xcolor}
\usepackage[hyphens]{url}
\usepackage[hidelinks]{hyperref}

\newcommand{\iscasubmissionnumber}{NaN}
\newcommand{\mechanism}{SAPA}
\begin{document}
\setlength{\paperwidth}{8.5in}
\setlength{\paperheight}{11in}
\ifdefined\pdfpagewidth
  \pdfpagewidth=8.5in
  \pdfpageheight=11in
\else
  \special{papersize=8.5in,11in}
\fi
\pagenumbering{arabic}
\title{Beyond Capacity: Service-Aware Memory Pooling\\
for Ethernet Scale-Up Inference}
\author{\normalsize{ISCA 2026 Submission
  \textbf{\#\iscasubmissionnumber} -- Confidential Draft -- Do NOT Distribute!!}}
\maketitle
\thispagestyle{plain}
\pagestyle{plain}

\begin{abstract}
Remote memory can expand the working sets of inference systems whose
local high-bandwidth memory is insufficient. However, fitting a working
set does not establish that its accesses can complete within an inference
deadline. This research proposal investigates the gap between physical
capacity and serviceable capacity in Ethernet scale-up memory pools.
We hypothesize that finite memory-service concurrency, endpoint transaction
windows, shared network paths, and receiver consumption jointly constrain
that gap. We outline service-aware placement and admission, which assigns
objects to memory tiers and budgets outstanding remote transactions using
calibrated service envelopes. The proposed evaluation separates capacity
expansion from additional hardware resources and distinguishes
representative-stage measurements from full inference service metrics.
This draft contains a proposed mechanism and evaluation plan, not measured
performance improvements or a completed hardware validation.
\end{abstract}

\section{Introduction}
The useful outcome of memory expansion is additional work completed
within a time budget. A larger memory pool may accommodate more KV state
while its transaction machinery cannot sustain the associated access
rate. Conversely, data that occupies substantial space but is rarely
accessed may be a favorable candidate for remote placement. We ask:
\emph{when does Ethernet memory pooling increase the inference working set
that a fixed accelerator deployment can actually serve?}

We use MPU to denote an independently attached memory-service node.
The proposed MPU provides remote storage and access service; it does not
execute attention or expert computation. Our scope is a multi-die
accelerator system where memory traffic and inference communication may
share the Ethernet scale-up fabric.

Several access patterns make nominal capacity an incomplete metric.
Active KV reads depend on the attention algorithm and selected state.
Paused sessions or cached prefixes can remain idle and later trigger
bursty restoration. Fine-grained tables can require many transactions
despite modest byte volume. These patterns must be distinguished:
older KV entries in dense attention are not automatically cold, and
restoring a paused session is not equivalent to reading its state on
every decode step.

\subsection{Research Hypothesis}
We hypothesize that capacity-only or nominal-bandwidth placement can
overestimate admissible load when memory-service slots, endpoint windows,
shared paths, or response consumption become limiting resources. A
placement and admission policy that accounts for these constraints may
increase the workload region satisfying an explicit time budget. The
hypothesis can fail: if calibrated workloads are primarily bandwidth
limited, a simple bandwidth-proportional policy may suffice.

The intended contributions are (1) a characterization of serviceable
capacity under specified access patterns, (2) an implementable placement
and transaction-admission mechanism, and (3) a calibrated evaluation of
its benefits and failure regions. None is claimed as completed in this
draft. In particular, adding more MPU nodes is not by itself evidence
for the mechanism: it also adds ports, service bandwidth, and slots.

\subsection{Positioning Against Existing Work}
Mooncake already combines a disaggregated KV cache with SLO-oriented
scheduling and prediction-based early rejection~\cite{mooncake}.
MemServe combines elastic memory pooling, context caching, and
locality-aware scheduling~\cite{memserve}. Fang and colleagues formulate
dynamic KV placement under capacity constraints to exploit heterogeneous
memory bandwidth~\cite{dynamicplacement}. Consequently, pooling,
capacity-aware placement, and admission control are not sufficient
novelty claims. The proposed distinction is the translation of concrete
Ethernet MPU transaction and response-path constraints into placement
and admission decisions. A full related-work review and experimental
comparison are still required to establish that distinction.

\section{From Physical to Serviceable Capacity}
\subsection{Capacity and Access Demand}
Let $x_{jp}$ denote bytes of class $j$ resident in pool $p$, and $C_p$
its usable capacity. For a specified observation period $\Delta$, let
$a_j$ be the fraction of resident bytes accessed and $A_{jp}$ the
applicable transfer amplification. A simplified demand model is
\begin{align}
  \sum_j x_{jp} &\le C_p, \\
  D_p = \frac{\sum_j x_{jp}a_jA_{jp}}{\Delta}
       &\le B_{p,\mathrm{useful}}.
  \label{eq:demand}
\end{align}
All terms must use consistent byte definitions. Read requests, responses,
writes, and on-wire overhead require separate resource accounting.
Constraints also apply to each die's HBM and each shared path; multiple
pools cannot each receive the full budget of the same bottleneck link.
Equation~\ref{eq:demand} is a necessary average-rate condition, not a
sufficient deadline or tail-latency guarantee.

\subsection{Distinct Concurrency Lifetimes}
An internal MPU slot is occupied from memory admission until data service
finishes. An endpoint transaction can remain outstanding until the
response reaches its specified acceptance or consumption point. Network
credits have their own allocation and return events. These lifetimes
must not be conflated.

For homogeneous transactions of $g$ useful bytes, a diagnostic upper
bound is
\begin{equation}
\begin{split}
 B_{p,\mathrm{useful}} \le \min\{&B_{\mathrm{mem}}, B_{\mathrm{path}},
 B_{\mathrm{recv}},\\
 &N_s g/L_s,\; N_e g/R_e,\; P u\},
\end{split}
\label{eq:bound}
\end{equation}
where $N_s$ and $L_s$ describe service slots and their mean occupancy,
$N_e$ and $R_e$ describe endpoint transactions, and $P$ and $u$ describe
packet throughput and useful bytes per packet. Terms must be converted
to useful-byte rates. The bound is a diagnostic approximation rather
than a new performance law; queueing, dependencies, and credit return
require explicit timing models.

\begin{figure}[t]
\centering
\setlength{\fboxsep}{6pt}
\fbox{\begin{minipage}{0.91\columnwidth}
\centering
Per-die placement and remote admission\\[3pt]
$\downarrow$ request path / shared fabric\\[3pt]
MPU admission $\rightarrow$ latency $\rightarrow$ data service\\
\emph{Internal service slot released here}\\[3pt]
$\downarrow$ response queue / shared fabric\\[3pt]
Finite receiver buffer $\rightarrow$ HBM consumption\\
\emph{Proposed endpoint budget released at consumption}
\end{minipage}}
\caption{Proposed control boundary. Internal memory slots, network
credits, and endpoint transaction budgets have separate lifetimes.
This is a mechanism schematic, not a measured execution trace.}
\label{fig:lifetime}
\end{figure}

An illustrative calculation uses the current configuration's 128 service
slots, 256-byte requests, and 200-ns fixed memory latency:
\begin{equation}
  \frac{128\times256\;\mathrm{B}}{200\;\mathrm{ns}}
  = 163.84\;\mathrm{GB/s}.
\end{equation}
This loose upper bound is below the configured 190-GB/s memory bandwidth
even before data service and queueing extend occupancy. It is arithmetic
from assumed parameters, not a hardware measurement. Approximately 149
slots cover only the fixed latency at that rate; they do not guarantee
full throughput. Substituting network round-trip time into this internal
slot calculation would misrepresent the current model.

\subsection{Metric Definition}
For a fixed access class, time budget, and hardware allocation, we define
\emph{serviceable capacity} as the largest resident working set satisfying
the specified completion budget. Mixed classes produce a feasible region,
not a universal scalar capacity. Ultimately, deployment value should be
measured by completed requests meeting both TTFT and TPOT objectives
under a fixed external arrival stream, with rejected work reported.
Representative-stage simulation alone cannot establish those metrics.

\section{Proposed Placement and Admission}
We provisionally call the mechanism Service-Aware Placement and Admission
(\mechanism). The initial implementation would combine an offline
placement planner with finite endpoint transaction budgets, keeping
packet packing and VC policies fixed.

\subsection{Planner Inputs and Decisions}
Inputs include object sizes, observed access counts, stage dependencies,
per-die free HBM, Host and MPU capacity, calibrated service curves, and
the mapping of traffic onto shared paths. Decisions assign objects or
shards to memory tiers and choose endpoint outstanding windows for
traffic classes. The first version avoids replication to keep capacity
accounting and consistency manageable.

A practical greedy planner can evaluate candidate migrations by the HBM
space released and their incremental demand on each constrained resource.
It prioritizes feasible migrations by released capacity per unit of
dominant-resource pressure, handling zero incremental demand separately.
After each assignment it recomputes shared-path demand. Hot objects are
sharded only when the resulting service parallelism outweighs additional
transactions and path overhead. The objective is to enlarge the admitted
working set under the budget, rather than obtain low latency merely by
serving fewer requests. Small-instance exhaustive search provides a
reference for the quality of this heuristic.

\subsection{Transaction Admission Boundary}
The endpoint acquires an outstanding-transaction token before issuing
remote work. Response storage must be finite and accounted for; the
proposed token is released after the associated data reaches the chosen
HBM-consumption point, as illustrated in Fig.~\ref{fig:lifetime}.
Internal MPU slots remain separate and can be released earlier.
Admission tokens are not CBFC credits and do not change protocol credit
semantics, return-port mapping, or transaction ordering.

The implementation must specify response reservations, bounded waiting,
and resource release on completion or failure. New requests must not
consume all resources needed by responses. Progress tests are necessary;
a smaller window alone is not a deadlock argument. Planner overhead,
state size, and sensitivity to inaccurate access estimates belong in
the evaluation.

\subsection{Scope of the First Version}
Service curves are fitted using calibration workloads, then frozen
before evaluating unseen traces. A dynamic controller is unnecessary
until static configurations establish useful headroom. If a workload
cannot satisfy the budget, the policy may choose a different tier or
reject additional work; both rejection and unserved demand are visible
outcomes. Whole-system service guarantees require an arrival and batching
model beyond the current representative-stage simulator.

\section{Evaluation Plan and Evidence Status}
\subsection{Available Model and Missing Mechanisms}
The existing platform provides capacity accounting, representative
inference DAGs, AXI packetization, finite MPU memory-service concurrency,
Ethernet credit flow control, shared HBM service, and a transaction-level
PCIe comparison path. These are behavioral transaction and packet models;
picosecond timestamp resolution does not make them instruction-level or
cycle-accurate GPU simulation.

Several gaps matter directly to the proposed claims. The current MPU
slot is released after memory service. Pending admission and pending
responses do not yet constitute a demonstrated complete finite-buffer
feedback chain. SUE data is accumulated at message granularity before
HBM writing, rather than using a complete packet-level credit/consumer
coupling. Aggregate issue limits must be defined per port or die, since
a per-task 1-ns issue interval does not establish a hardware issue rate.
Packet-rate counters do not impose a switch packet-processing limit.
If packet throughput is central to a result, the switch model needs
separate pipeline latency, initiation interval, parallel capacity, and
shared arbitration.

Current traces use synthesized routing and addresses, and do not
establish production locality. Object-level HBM residency, online
migration, continuous batching, and complete attention execution are not
assumed to be available. The proposed mechanism is not implemented in
this manuscript project.

\subsection{Calibration and Mechanism Experiments}
First, calibrate the model using measurements from real endpoints or
FPGA/RTL prototypes. Sweep request size, concurrency, read/write mix,
and response consumption rate, recording latency and throughput.
Separate fitting data from held-out verification cases. Measure both
service-slot occupancy and end-to-end occupancy, checking payload and
wire-byte conservation. If only individual subsystems can be calibrated,
state the resulting limit on system-level confidence.

Next, hold capacity and bandwidth fixed while varying service slots,
endpoint windows, receiver buffers, locality, and shared MoE traffic.
Access classes include active KV, session restoration, and fine-grained
lookup where supported by validated traces. This stage asks which
resource limits serviceable capacity, without claiming production
inference throughput.

\subsection{Baselines and Resource Fairness}
Baselines should include HBM-only with feasible batches, Host memory via
PCIe, all-remote MPU placement, uniform sharding, capacity-first,
bandwidth-proportional, and locality-based placement. Compare these with
\mechanism\ and small-instance static optima. Preserve accelerator count,
per-die HBM, remote capacity, links, service bandwidth, total slots,
buffers, and credits across policy comparisons.

Node-count scaling is a separate hardware-expansion experiment with an
explicit resource table. An HBM-only configuration that cannot admit a
large batch is a capacity rejection, not an infinite-latency performance
point. It should also be allowed to process smaller batches under the
same external workload. Cost or power claims require their own validated
models and are outside the initial scope.

\subsection{Ablations and Decisive Figures}
Model ablations progressively add memory-service concurrency, endpoint
windows, finite response storage, and shared resource timing to a
capacity/bandwidth baseline. Policy ablations remove placement or
transaction admission independently. Use unseen access patterns and
load levels after selecting parameters.

The primary proposed figure compares the physical-capacity frontier,
bandwidth-only prediction, detailed-model feasible region, and the
region attained by \mechanism. A second figure tests whether model
detail reverses a placement decision under identical hardware. Such a
reversal is a hypothesis, not a required outcome; numerical rounding or
unequal resources must not manufacture it. A third figure traces service
slots, response buffers, and endpoint budgets over time to explain any
observed boundary.

\subsection{Reporting and Falsification}
Initially report stage completion, transaction latency, useful and wire
bytes, occupancy, actual waiting time, and admitted working set. Summed
per-request waiting time is not a critical-path duration. Capacity
rejections, timeouts, and incomplete runs remain distinct from valid
performance points. End-to-end request goodput requires complete
inference dependencies, attention reads/writes, arrivals, and batching;
multiplying a representative block by layer count is insufficient.

The central architectural claim should be reduced or abandoned if
held-out workloads are adequately explained by bandwidth alone, if simple
baselines match the proposed policy, or if realistic finite buffering
eliminates the expected advantage. This draft provides no performance
results, new simulator executions, or hardware calibration evidence.

\bibliographystyle{IEEEtranS}
\bibliography{refs}
\end{document}
